\documentclass[a4paper,11pt]{article}

% =========================
% Pacotes básicos
% =========================
\usepackage[utf8]{inputenc}
\usepackage[T1]{fontenc}
\usepackage[brazil]{babel}
\usepackage{graphicx}
\usepackage{hyperref}
\usepackage{amsmath}
\usepackage{amssymb}
\usepackage[margin=2.0cm]{geometry}
\usepackage{booktabs}
\usepackage{tabularx}
\usepackage{array}
\usepackage{xcolor}
\usepackage{listings}
\usepackage{float}
\usepackage{enumitem}
\usepackage{caption}

% =========================
% Configurações de layout
% =========================
\setlength{\parskip}{0.2em}
\setlength{\parindent}{0pt}
\setlength{\textfloatsep}{8pt plus 1pt minus 2pt}
\setlength{\floatsep}{6pt plus 1pt minus 2pt}
\setlength{\intextsep}{8pt plus 1pt minus 2pt}
\setlength{\abovecaptionskip}{4pt}
\setlength{\belowcaptionskip}{2pt}
\setlist[itemize]{noitemsep, topsep=2pt, leftmargin=1.5em}

% =========================
% Configuração para blocos de código
% =========================
\lstset{
  basicstyle=\ttfamily\scriptsize,
  breaklines=true,
  frame=single,
  columns=fullflexible,
  keepspaces=true,
  showstringspaces=false
}

\graphicspath{{reports_iqn_temporal/}{./}}

% =========================
% Figuras opcionais para Overleaf
% =========================
% Se voce colar apenas este .tex no Overleaf, o documento compila com caixas
% indicando as figuras ausentes. Para ver as figuras reais, envie tambem a pasta
% reports_iqn_temporal/ com os arquivos .png gerados.
\newcommand{\imagemsegura}[2][]{%
  \IfFileExists{#2}{%
    \includegraphics[#1]{#2}%
  }{%
    \fbox{%
      \begin{minipage}[c][4cm][c]{0.9\linewidth}
      \centering\scriptsize
      Figura nao encontrada no Overleaf.\\
      Envie o arquivo: \texttt{#2}
      \end{minipage}%
    }%
  }%
}

% =========================
% Dados editáveis do trabalho
% =========================
\newcommand{\tituloTrabalho}{Previsão de Séries Temporais de SSH com GRU, Codificação Temporal e IQN}
\newcommand{\nomeAutor}{Matheus Costa Monteiro}
\newcommand{\numeroUSP}{17756099}
\newcommand{\disciplina}{PCS5024 -- Aprendizado Estatístico}
\newcommand{\dataEntrega}{\today}
\newcommand{\linkRepositorio}{}

\begin{document}

% =========================
% Cabeçalho
% =========================
\begin{center}
\begin{minipage}{0.48\textwidth}
    \raggedright
    \IfFileExists{pcs.png}{\includegraphics[height=1.1cm]{pcs.png}}{\fbox{\scriptsize PCS}}
\end{minipage}
\begin{minipage}{0.48\textwidth}
    \raggedleft
    \IfFileExists{poli.png}{\includegraphics[height=1.6cm]{poli.png}}{\fbox{\scriptsize POLI-USP}}
\end{minipage}

\vspace{0.4em}

{\large\bfseries \tituloTrabalho\par}
\vspace{0.2em}
{\normalsize \nomeAutor{} -- \numeroUSP\par}
{\small \disciplina\par}
{\small \dataEntrega\par}
\end{center}

% =========================
% Resumo
% =========================
\begin{abstract}\small
Este trabalho avalia modelos recorrentes para previsão da série temporal de nível
do mar SSH (\textit{sea surface height}) sob diferentes níveis de dados faltantes.
Foram comparadas duas arquiteturas determinísticas baseadas em GRU: uma versão
base, sem informação temporal explícita no decodificador, e uma versão com
codificação temporal sinusoidal inspirada em Vaswani et al. \cite{vaswani2017}.
Além disso, foi implementada uma extensão probabilística com IQN
(\textit{Implicit Quantile Networks}), conforme Gouttes et al. \cite{gouttes2021},
para estimar quantis e intervalos de previsão. Os experimentos foram executados
com razões de remoção de dados iguais a $0{,}0$, $0{,}1$, $0{,}3$ e $0{,}5$,
totalizando doze cenários. A codificação temporal reduziu o RMSE em todos os
níveis de dados faltantes, com ganhos relativos de aproximadamente $22\%$ a
$41\%$ em relação ao GRU base. A IQN produziu previsões probabilísticas com
cobertura empírica entre $71\%$ e $79\%$ para o intervalo nominal de $80\%$ e
entre $84\%$ e $89\%$ para o intervalo nominal de $90\%$.
\end{abstract}

\section{Objetivo e escopo}

O objetivo do experimento foi comparar a robustez de modelos GRU para previsão
de SSH quando parte das observações da série temporal é removida e, em seguida,
estender o pipeline temporal para produzir previsões probabilísticas. A
comparação determinística foi organizada em oito cenários, quatro para o GRU
base e quatro para o GRU com codificação temporal. A etapa probabilística
adicionou quatro cenários com IQN acoplada ao GRU temporal, usando os mesmos
níveis de \textit{missing}. Assim, o relatório cobre tanto previsão pontual
quanto estimativas de incerteza por quantis.

\section{Base de dados e preparação}

A base utilizada foi o arquivo \texttt{data/santos\_ssh.csv}, contendo
observações de SSH associadas a timestamps. O conjunto foi dividido por data:
observações anteriores a \texttt{2020-06-01 00:00:00} foram usadas para treino,
e observações a partir dessa data foram usadas para teste. A normalização foi
calculada apenas sobre o conjunto de treino e depois aplicada aos conjuntos de
treino e teste, evitando vazamento de informação.

As amostras de entrada foram criadas por janelas deslizantes. Cada janela utiliza
um contexto passado de até $8000$ minutos, correspondente a aproximadamente
$800$ observações quando a série está completa, e prevê um horizonte futuro de
até $2000$ minutos, aproximadamente $200$ observações. A cada nível de
\textit{missing}, uma fração das linhas do treino e do teste foi removida de
forma aleatória e reprodutível antes da formação das janelas. Por isso, nos
gráficos de previsão o eixo horizontal indica o passo futuro válido da janela, e
não o número de épocas de treinamento.

\section{Modelos}

\subsection{GRU base}

O modelo base utiliza uma GRU no codificador para resumir a sequência passada de
SSH. O decodificador também é uma GRU, mas recebe entradas nulas durante a
previsão, de forma que a informação disponível para gerar o futuro vem
essencialmente do estado oculto final produzido pelo codificador. Esse modelo
serve como referência para medir quanto a codificação temporal explícita melhora
a capacidade de previsão.

\subsection{GRU com codificação temporal}

O modelo temporal adiciona uma codificação sinusoidal dos timestamps relativos
de cada janela. A codificação segue a ideia das representações posicionais de
Vaswani et al. \cite{vaswani2017}, com componentes senoidais e cossenoidais em
diferentes escalas:

\begin{equation}
    PE(t, 2i) = \sin\left(\frac{t}{10000^{2i/d}}\right),
    \qquad
    PE(t, 2i+1) = \cos\left(\frac{t}{10000^{2i/d}}\right).
\end{equation}

No codificador, essa representação é concatenada à feature SSH, produzindo uma
entrada do tipo $[SSH, PE(t)]$. No decodificador, a GRU recebe a codificação
temporal dos instantes futuros, permitindo que o modelo use explicitamente a
posição relativa de cada passo previsto. Os timestamps são sempre relativos à
origem da janela de previsão, evitando o uso direto de timestamps absolutos.

\subsection{GRU temporal com IQN}

A extensão probabilística reutiliza o codificador e o decodificador do GRU
temporal, mas substitui a cabeça determinística por uma cabeça IQN. Durante o
treinamento, para cada observação válida e passo futuro é amostrado um valor
$\tau \sim U(0,1)$. Esse valor é codificado por uma base cossenoidal,
seguindo a formulação de Gouttes et al. \cite{gouttes2021}:

\begin{equation}
    \phi(\tau) = ReLU\left(\sum_{i=0}^{n-1}
    \cos(\pi i \tau)w_i + b_i\right).
\end{equation}

O vetor $\phi(\tau)$ é combinado ao estado oculto produzido pelo decodificador
por produto elemento a elemento na forma $h_t \odot (1 + \phi(\tau))$. Em
seguida, uma pequena rede MLP gera a previsão associada ao quantil $\tau$. Na
inferência, foram usadas $100$ amostras determinísticas de $\tau$ para estimar
os quantis empíricos $q05$, $q10$, $q50$, $q90$ e $q95$. A mediana $q50$ foi
usada como previsão pontual da IQN.

\section{Treinamento e métricas}

Todos os cenários foram treinados com os mesmos hiperparâmetros principais:
\begin{itemize}
    \item limite máximo de $1000$ épocas;
    \item regularização por \textit{early stopping}, com paciência de $50$
    épocas sem melhora na loss de teste;
    \item GRU com dimensão oculta $128$;
    \item dimensão da codificação temporal $128$;
    \item batch size efetivo $1024$, obtido com batch $256$ por processo;
    \item taxa de aprendizado $2 \times 10^{-4}$;
    \item decaimento de peso $10^{-5}$;
    \item semente aleatória $100$.
\end{itemize}

Nos cenários IQN, foram usadas $8$ amostras de $\tau$ por ponto futuro durante
o treinamento, $100$ amostras durante a avaliação e dimensão $64$ para o
embedding de $\tau$.

As execuções foram realizadas em treinamento distribuído com quatro GPUs RTX
4000 Ada Generation, usando um processo por GPU. Ao final do treinamento, o
checkpoint da melhor época foi restaurado antes da geração das previsões e do
cálculo das métricas finais.

Nos modelos determinísticos, a função objetivo foi o erro quadrático médio
calculado apenas sobre posições válidas das sequências, ignorando o padding
criado pela formação das janelas. Nos cenários IQN, a função objetivo foi a
\textit{pinball loss}, também chamada de \textit{quantile loss}:

\begin{equation}
    L_\tau(y, \hat{y}) =
    \max \left(\tau (y-\hat{y}), (\tau-1)(y-\hat{y})\right).
\end{equation}

Assim como no MSE, a loss de quantis foi calculada apenas nas posições válidas,
desconsiderando o padding. As métricas finais foram calculadas na escala
original da SSH:

\begin{equation}
    MSE = \frac{1}{N}\sum_{i=1}^{N}(y_i - \hat{y}_i)^2,
    \qquad
    RMSE = \sqrt{MSE},
    \qquad
    MAE = \frac{1}{N}\sum_{i=1}^{N}|y_i - \hat{y}_i|.
\end{equation}

Para IQN, as métricas pontuais usam $q50$ como $\hat{y}$. Também foram
calculadas a quantile loss, a cobertura empírica dos intervalos $q10$--$q90$ e
$q05$--$q95$, e a largura média desses intervalos. A coluna \textit{Loss teste}
vem da função objetivo usada durante o treinamento, na escala normalizada. Já a
coluna QLoss reportada nas tabelas foi recalculada sobre os quantis
desnormalizados, na escala original da SSH.

\section{Implementação no pipeline}

A implementação foi concentrada em dois pontos do código. Em
\texttt{PCS5024\_EP\_Time\_Series.py}, foi adicionada a variante
\texttt{gru\_temporal\_iqn}. Essa variante reaproveita o codificador e o
decodificador do GRU temporal e troca apenas a cabeça de saída por uma cabeça
IQN. O treinamento amostra valores de $\tau$ a cada passo válido da sequência,
aplica a codificação cossenoidal de $\tau$, combina essa representação ao estado
oculto do decodificador e otimiza a \textit{pinball loss} mascarada por
\texttt{target\_lengths}. Na avaliação, uma grade determinística de $100$
valores de $\tau$ é usada para estimar os quantis exportados.

Para manter compatibilidade com o restante do pipeline, a mediana $q50$ é
gravada também como \texttt{Y\_hat}. Assim, os CSVs e as métricas pontuais
continuam tendo o mesmo formato dos cenários determinísticos, enquanto os
cenários IQN adicionam as colunas \texttt{q05}, \texttt{q10}, \texttt{q50},
\texttt{q90}, \texttt{q95}, \texttt{interval\_width\_80} e
\texttt{interval\_width\_90}. Os arquivos JSON de métricas também passaram a
incluir \texttt{quantile\_loss}, coberturas empíricas e larguras médias dos
intervalos.

Em \texttt{run\_reports.py}, foram preservados os cenários determinísticos A--H
e adicionados os cenários I--L com \texttt{gru\_temporal\_iqn} para
\textit{missing} $0{,}0$, $0{,}1$, $0{,}3$ e $0{,}5$. O gerador de relatórios
também foi atualizado para produzir gráficos com bandas probabilísticas nos
cenários IQN, além dos gráficos agregados de cobertura nominal versus empírica e
largura média dos intervalos. A execução final foi salva na pasta
\texttt{reports\_iqn\_temporal}.

\section{Cenários experimentais}

Foram executados doze cenários, resumidos na Tabela~\ref{tab:resultados}. Os
cenários A--D usam o GRU base, E--H usam o GRU temporal determinístico, e I--L
usam o GRU temporal com IQN. As razões de dados faltantes foram $0{,}0$,
$0{,}1$, $0{,}3$ e $0{,}5$.

\begin{table}[H]
\centering
\caption{Resultados finais dos doze cenários com \textit{early stopping}. A
loss de teste é MSE normalizado para modelos determinísticos e quantile loss
normalizada para IQN. MSE, RMSE e MAE estão na escala original da SSH; nos
cenários IQN, essas métricas usam a mediana $q50$. A coluna QLoss é calculada na
escala original da SSH.}
\label{tab:resultados}
\resizebox{\textwidth}{!}{%
\begin{tabular}{lllllrrrrrrr}
\hline
Cenario & Modelo & Missing & Epocas & Objetivo & Loss teste & MSE & RMSE & MAE & QLoss & Cob80 & Cob90 \\
\hline
A & GRU base & 0.0 & 337 & baseline ideal & 0.3565 & 0.0583 & 0.2415 & 0.1970 &  &  &  \\
B & GRU base & 0.1 & 196 & robustez & 0.4006 & 0.0660 & 0.2569 & 0.2064 &  &  &  \\
C & GRU base & 0.3 & 140 & degradação principal & 0.5017 & 0.0834 & 0.2888 & 0.2295 &  &  &  \\
D & GRU base & 0.5 & 334 & degradação severa & 0.4769 & 0.0696 & 0.2637 & 0.2120 &  &  &  \\
E & GRU temporal & 0.0 & 340 & controle & 0.2150 & 0.0355 & 0.1884 & 0.1496 &  &  &  \\
F & GRU temporal & 0.1 & 756 & robustez & 0.1902 & 0.0312 & 0.1767 & 0.1389 &  &  &  \\
G & GRU temporal & 0.3 & 664 & recuperação principal & 0.1750 & 0.0288 & 0.1697 & 0.1328 &  &  &  \\
H & GRU temporal & 0.5 & 630 & robustez severa & 0.2110 & 0.0331 & 0.1818 & 0.1445 &  &  &  \\
I & GRU temporal IQN & 0.0 & 869 & IQN controle & 0.1093 & 0.0259 & 0.1611 & 0.1250 & 0.0304 & 0.7680 & 0.8759 \\
J & GRU temporal IQN & 0.1 & 976 & IQN robustez & 0.1095 & 0.0249 & 0.1577 & 0.1240 & 0.0299 & 0.7283 & 0.8564 \\
K & GRU temporal IQN & 0.3 & 342 & IQN recuperação principal & 0.1408 & 0.0423 & 0.2058 & 0.1634 & 0.0387 & 0.7867 & 0.8936 \\
L & GRU temporal IQN & 0.5 & 590 & IQN robustez severa & 0.1292 & 0.0324 & 0.1800 & 0.1426 & 0.0343 & 0.7136 & 0.8396 \\
\hline
\end{tabular}%
}
\end{table}

\section{Resultados determinísticos}

A Figura~\ref{fig:comparativo} mostra a evolução do RMSE e do MAE em função da
taxa de dados faltantes. O GRU temporal manteve erro menor que o GRU base em
todos os níveis de \textit{missing}. Em RMSE, as reduções relativas foram
aproximadamente $22\%$ para \textit{missing} $0{,}0$, $31\%$ para $0{,}1$,
$41\%$ para $0{,}3$ e $31\%$ para $0{,}5$. Em MAE, os ganhos ficaram entre
$24\%$ e $42\%$.

\begin{figure}[H]
    \centering
    \imagemsegura[width=0.92\textwidth]{reports_iqn_temporal/comparativo_rmse_mae.png}
    \caption{Comparação agregada de RMSE e MAE por taxa de dados faltantes.}
    \label{fig:comparativo}
\end{figure}

Esse resultado confirma que a codificação temporal ajuda a GRU a interpretar
janelas irregulares. No modelo base, remover linhas altera os intervalos reais
entre observações, mas o decodificador não recebe essa informação. No modelo
temporal, tanto o contexto passado quanto o horizonte futuro carregam posições
relativas, o que reduz a degradação quando há dados faltantes.

\section{Resultados probabilísticos da IQN}

A IQN foi avaliada nos cenários I--L. A mediana $q50$ foi usada como previsão
pontual e os intervalos $q10$--$q90$ e $q05$--$q95$ foram usados para avaliar
cobertura. A Tabela~\ref{tab:iqn} resume as métricas probabilísticas e as
larguras médias dos intervalos.

\begin{table}[H]
\centering
\caption{Métricas probabilísticas dos cenários IQN.}
\label{tab:iqn}
\resizebox{\textwidth}{!}{%
\begin{tabular}{lrrrrrrr}
\hline
Cenario & Missing & QLoss & RMSE q50 & MAE q50 & Cob80 & Cob90 & Largura 80/90 \\
\hline
I & 0.0 & 0.0304 & 0.1611 & 0.1250 & 0.7680 & 0.8759 & 0.3872 / 0.5067 \\
J & 0.1 & 0.0299 & 0.1577 & 0.1240 & 0.7283 & 0.8564 & 0.3474 / 0.4590 \\
K & 0.3 & 0.0387 & 0.2058 & 0.1634 & 0.7867 & 0.8936 & 0.5111 / 0.6648 \\
L & 0.5 & 0.0343 & 0.1800 & 0.1426 & 0.7136 & 0.8396 & 0.3876 / 0.5028 \\
\hline
\end{tabular}%
}
\end{table}

\begin{figure}[H]
\centering
\begin{minipage}{0.49\textwidth}
    \centering
    \imagemsegura[width=\textwidth]{reports_iqn_temporal/cobertura_iqn.png}
    \caption*{Cobertura nominal versus cobertura empírica.}
\end{minipage}
\begin{minipage}{0.49\textwidth}
    \centering
    \imagemsegura[width=\textwidth]{reports_iqn_temporal/largura_intervalos_iqn.png}
    \caption*{Largura média dos intervalos por \textit{missing}.}
\end{minipage}
\caption{Avaliação dos intervalos probabilísticos produzidos pela IQN.}
\label{fig:iqn-cobertura}
\end{figure}

As coberturas ficaram ligeiramente abaixo dos valores nominais na maioria dos
cenários, indicando intervalos um pouco estreitos. Mesmo assim, os resultados
mostram calibração razoável: o intervalo nominal de $80\%$ cobriu entre
$71\%$ e $79\%$ dos pontos, enquanto o intervalo nominal de $90\%$ cobriu entre
$84\%$ e $89\%$. A maior largura ocorreu no cenário K, com \textit{missing}
$0{,}3$, que também foi o cenário IQN com maior RMSE da mediana. Isso sugere
que a IQN aumentou a incerteza justamente em uma configuração mais difícil.

Comparando a mediana da IQN ao GRU temporal determinístico, a IQN reduziu o RMSE
nos cenários sem remoção, com \textit{missing} $0{,}1$ e com \textit{missing}
$0{,}5$. No cenário com \textit{missing} $0{,}3$, a mediana da IQN teve RMSE
maior que o GRU temporal determinístico, mas ainda produziu intervalos de
incerteza com cobertura próxima dos valores nominais.

\section{Curvas de treinamento}

As curvas de loss indicam que os cenários atingem seus melhores checkpoints em
épocas distintas. Por isso, o treinamento foi interrompido quando a loss de teste
ficou $50$ épocas sem melhora. Na IQN, as curvas representam quantile loss, e
não MSE, o que torna a comparação direta de loss com os cenários determinísticos
menos apropriada.

\begin{figure}[H]
\centering
\begin{minipage}{0.49\textwidth}
    \centering
    \imagemsegura[width=\textwidth]{reports_iqn_temporal/cenario_B_GRU_base_missing_0p1_loss.png}
    \caption*{Cenário B: GRU base, missing $0{,}1$.}
\end{minipage}
\begin{minipage}{0.49\textwidth}
    \centering
    \imagemsegura[width=\textwidth]{reports_iqn_temporal/cenario_F_GRU_temporal_missing_0p1_loss.png}
    \caption*{Cenário F: GRU temporal, missing $0{,}1$.}
\end{minipage}

\vspace{0.5em}

\begin{minipage}{0.49\textwidth}
    \centering
    \imagemsegura[width=\textwidth]{reports_iqn_temporal/cenario_J_GRU_temporal_IQN_missing_0p1_loss.png}
    \caption*{Cenário J: GRU temporal IQN, missing $0{,}1$.}
\end{minipage}
\begin{minipage}{0.49\textwidth}
    \centering
    \imagemsegura[width=\textwidth]{reports_iqn_temporal/cenario_K_GRU_temporal_IQN_missing_0p3_loss.png}
    \caption*{Cenário K: GRU temporal IQN, missing $0{,}3$.}
\end{minipage}
\caption{Curvas de loss de treino e teste para cenários representativos.}
\label{fig:losses}
\end{figure}

\section{Exemplos de previsão}

A Figura~\ref{fig:forecasts} mostra exemplos de previsão em janelas de teste
representativas. Nos cenários determinísticos, o gráfico mostra a previsão
pontual. Nos cenários IQN, o gráfico mostra a mediana $q50$ e bandas de
incerteza $q10$--$q90$ e $q05$--$q95$.

\begin{figure}[H]
\centering
\begin{minipage}{0.49\textwidth}
    \centering
    \imagemsegura[width=\textwidth]{reports_iqn_temporal/cenario_B_GRU_base_missing_0p1_forecast.png}
    \caption*{Cenário B: GRU base, missing $0{,}1$.}
\end{minipage}
\begin{minipage}{0.49\textwidth}
    \centering
    \imagemsegura[width=\textwidth]{reports_iqn_temporal/cenario_F_GRU_temporal_missing_0p1_forecast.png}
    \caption*{Cenário F: GRU temporal, missing $0{,}1$.}
\end{minipage}

\vspace{0.5em}

\begin{minipage}{0.49\textwidth}
    \centering
    \imagemsegura[width=\textwidth]{reports_iqn_temporal/cenario_J_GRU_temporal_IQN_missing_0p1_forecast.png}
    \caption*{Cenário J: IQN, missing $0{,}1$.}
\end{minipage}
\begin{minipage}{0.49\textwidth}
    \centering
    \imagemsegura[width=\textwidth]{reports_iqn_temporal/cenario_L_GRU_temporal_IQN_missing_0p5_forecast.png}
    \caption*{Cenário L: IQN, missing $0{,}5$.}
\end{minipage}
\caption{Exemplos de previsão no conjunto de teste.}
\label{fig:forecasts}
\end{figure}

\section{Discussão}

Os resultados mostram que a codificação temporal não apenas melhora o erro em
dados completos, mas também reduz a degradação quando há dados faltantes. A
principal diferença arquitetural é que o modelo temporal recebe informação sobre
o instante relativo de cada observação passada e de cada passo futuro. Quando a
série fica irregular por remoção de linhas, essa informação ajuda a GRU a
interpretar intervalos temporais que deixam de ser uniformes.

A implementação da IQN complementou o pipeline ao transformar a previsão pontual
em previsão probabilística. A maior dificuldade prática foi que a cabeça IQN
multiplica o custo computacional: no treino são geradas $8$ amostras por ponto
futuro e, na avaliação, $100$ amostras por ponto futuro. Além disso, no
treinamento distribuído foi necessário remover a cabeça linear determinística
dos modelos IQN para evitar parâmetros não utilizados no DDP. Após esse ajuste,
os cenários I--L executaram corretamente.

Uma limitação importante do protocolo é que o \textit{early stopping} monitorou
a loss do conjunto de teste, pois este era o particionamento disponível no
pipeline experimental. Em um protocolo mais rigoroso, uma fração do treino
deveria ser separada como validação, reservando o teste apenas para a avaliação
final. Outra limitação é que as coberturas da IQN ficaram levemente abaixo dos
níveis nominais, sugerindo que uma calibração posterior ou um ajuste de
hiperparâmetros poderia melhorar a largura dos intervalos.

\section{Conclusão}

A comparação experimental indica que o GRU com codificação temporal é superior
ao GRU base nos oito cenários determinísticos avaliados. O ganho é mais
expressivo nos cenários com maior proporção de dados faltantes: com
\textit{missing} $0{,}3$, o RMSE do GRU temporal foi $0{,}1697$, contra
$0{,}2888$ do GRU base; com \textit{missing} $0{,}5$, foi $0{,}1818$, contra
$0{,}2637$. Portanto, a concatenação de uma representação temporal sinusoidal ao
SSH, aliada ao uso de informação temporal futura no decodificador, mostrou-se
uma estratégia eficaz para previsão de SSH em séries irregulares.

A IQN atendeu ao objetivo probabilístico do trabalho, pois produziu quantis,
bandas de incerteza e métricas de cobertura. A mediana da IQN alcançou RMSE
competitivo e, em três dos quatro níveis de \textit{missing}, foi menor ou muito
próxima ao RMSE do GRU temporal determinístico. As coberturas empíricas ficaram
próximas, mas ligeiramente abaixo, das coberturas nominais, indicando que o
modelo captura a incerteza principal da série, embora ainda possa ser calibrado
para intervalos mais conservadores.

% =========================
% Referências
% =========================
{\footnotesize
\begin{thebibliography}{99}

\bibitem{vaswani2017}
VASWANI, A.; SHAZEER, N.; PARMAR, N.; USZKOREIT, J.; JONES, L.; GOMEZ, A. N.;
KAISER, L.; POLOSUKHIN, I. \textit{Attention Is All You Need}. Advances in
Neural Information Processing Systems, 2017. Disponível em:
\url{https://arxiv.org/abs/1706.03762}.

\bibitem{gouttes2021}
GOUTTES, A.; RASUL, K.; KOREN, M.; STEPHAN, J.; NAGHIBI, T.
\textit{Probabilistic Time Series Forecasting with Implicit Quantile Networks}.
Proceedings of the International Conference on Machine Learning, Time Series
Workshop, 2021. Disponível em: \url{https://arxiv.org/abs/2107.03743}.

\end{thebibliography}
}

\end{document}
